% Generated draft: requires mathematical and source review.
\documentclass[11pt]{article}
\usepackage[margin=25mm]{geometry}
\usepackage{fontspec}
\setmainfont{Latin Modern Roman}
\usepackage{amsmath,amssymb,mathtools,unicode-math}
\setmathfont{Latin Modern Math}
\usepackage{enumitem,needspace,xurl,hyperref}
\setlength{\parindent}{0pt}
\setlength{\parskip}{5pt}
\setlength{\emergencystretch}{3em}
\newcommand{\N}{\mathbb{N}}
\newcommand{\Z}{\mathbb{Z}}
\newcommand{\Q}{\mathbb{Q}}
\newcommand{\R}{\mathbb{R}}
\newcommand{\C}{\mathbb{C}}
\newcommand{\F}{\mathbb{F}}
\newcommand{\A}{\mathbb{A}}
\newcommand{\PP}{\mathbb{P}}
\newcommand{\E}{\mathbb{E}}
\newcommand{\eps}{\varepsilon}
\newcommand{\dd}{\,\mathrm{d}}
\begin{document}
\section*{UnsolvedMath ID: 6200015}
\textbf{Source problem code:} catalogue entry
\subsection{Definitions and mathematical statement}
Let $(N_1,\Delta_1)$ and $(N_2,\Delta_2)$ be closed $3$-manifolds with flag triangulations. A simplicial complex is flag if every finite set of pairwise adjacent vertices spans a simplex. The source writes $C(N_i,\Delta_i)$ for associated spaces and assumes
\[
 \partial_\infty C(N_1,\Delta_1)
       =\partial_\infty C(N_2,\Delta_2).
\]
Neither the construction $C$ nor its metric and boundary-at-infinity convention is defined in the record. The equality symbol could mean homeomorphism of topological boundaries or a more structured identification; that category must be fixed.

A nontrivial prime connected-sum factor of a $3$-manifold is a prime manifold occurring in a decomposition $N\cong P\mathbin\# Q$, with the sphere factor excluded in the intended convention. The question asks whether the boundary identification forces the same prime summand types in both directions. Formally, relative to the chosen manifold-isomorphism and prime conventions, the proposed conclusion is
\[
 \forall P\text{ prime},\quad
 \bigl(\exists Q_1:N_1\cong P\#Q_1\bigr)
 \ \Longleftrightarrow\
 \bigl(\exists Q_2:N_2\cong P\#Q_2\bigr).
\]
The original indexes the other manifold as $N_{i+1}$ for $i=1,2$; the displayed bidirectional reading treats the indices cyclically, so $N_3=N_1$. The record does not explicitly state that convention.

The wording asks whether each prime type appears in the other manifold, not explicitly whether the multiplicity of every prime factor agrees. Connectedness, orientability, orientation-preserving versus unoriented equivalence, and topological versus smooth category are not specified. The source additionally asks what can be said in higher dimensions; it gives no precise dimension range, analogue of the conclusion, or uniqueness assumptions for higher-dimensional decompositions.
\subsection{Short English statement}
If the spaces associated to two flag-triangulated closed three-manifolds have the same boundary at infinity, must the manifolds have the same prime connected-sum types? Ask also for appropriate higher-dimensional analogues. The construction of the associated spaces and the kind of boundary identification remain essential source dependencies.

The supplied summary describes the question as open and reports no resolution located. Its OPEN-TRIAGE label is not an independent verification. The 2005 source lists the problem, and provenance marks current status for review.
\subsection{Sources}
\begin{itemize}
\item Local UnsolvedMath registry, permanent ID 6200015, source code catalogue entry, ``Boundaries of Groups and Kleinian Groups --- Problem 15,'' proposed by Misha Kapovich and Tadeusz Januszkiewicz, 2005. Dataset curation and metadata: Ulam AI and UnsolvedMath Contributors, CC BY 4.0.
\item Kapovich, \emph{Problems on Boundaries of Groups and Kleinian Groups} (2005), Problem 15, PDF page 5, \url{https://www.math.ucdavis.edu/~kapovich/EPR/problems.pdf}; normalized pdf-text extraction, accessed 2026-07-29.
\item The supplied summary provides no later primary resolution reference or definition of $C(N,\Delta)$. Source-specific provenance marks current status and statement-reuse terms NEEDS\_REVIEW for the public author-hosted PDF, separately from the dataset licence.
\end{itemize}
\end{document}
